% ==================== 第 2 章 国内外研究现状 ====================
\clearpage
\section{国内外研究现状}

\subsection{排课与组合优化研究}

排课问题属于典型的组合优化问题。早期研究常将课程与时段冲突转化为图着色
或整数规划模型；随后，遗传算法、模拟退火、禁忌搜索等元启发式方法被用于
较大规模实例。此类方法能够在限定时间内搜索较优方案，但通常需要针对业务
设计修复算子，若处理不当，可能生成违反硬约束的中间解。

约束规划强调以变量、定义域和约束直接描述问题。现代 CP-SAT 求解器进一步
融合约束传播、SAT 冲突学习和整数优化技术，适合表达“恰好一次”“至多一次”
和资源互斥等离散规则。其优势在于：当求解器返回可行状态时，模型中声明的
硬约束均被满足；同时可通过目标函数优化偏好与资源利用。

\subsection{智能排课系统研究}

现有排课系统大体可分为三类。第一类是面向学校固定课表的专用系统，重点
处理课程、教师和教室的周期性安排；第二类是培训机构 SaaS，强调教务流程、
客户管理与通知；第三类是面向研究的算法原型，重点比较求解策略和标准数据集。

对个性化辅导场景而言，固定课表方法难以覆盖频繁调课，通用 SaaS 的内部
求解逻辑通常不透明，而研究原型又往往缺少可直接使用的业务界面。除此之外，
“求出一张表”并不等于完成教务决策：实际使用还要求解释偏好取舍、展示受
影响对象、保留历史版本并支持审批发布。

\subsection{现有研究的不足}

结合本课题应用场景，现有方法主要存在以下不足：

\begin{enumerate}
  \item 硬约束和软偏好混合表达，导致结果是否可执行不够清晰；
  \item 侧重首次排课，对教师请假后的最小影响重排研究不足；
  \item 评价指标只给出单一总分，难以解释学生、教师与资源之间的取舍；
  \item 算法原型与协同平台脱节，数据维护、方案审批和消息通知难以形成闭环；
  \item 缺少独立于求解器的结果复核，系统正确性过度依赖单一实现。
\end{enumerate}

\subsection{本课题的切入点}

本课题不追求发明新的通用求解器，而是研究如何把成熟的 CP-SAT 能力转化为
可解释、可校验、可集成的教务决策系统。具体切入点包括：建立面向个性化
辅导的结构化约束模型；设计多策略权重与分项评分；以变更数量为主要目标
实现局部重排；用独立校验器形成双重保障；最后把求解能力封装为 API，为
飞书多维表格、消息卡片和审批流程提供稳定接口。
